\subsection{The Even Root Property} 
With even powers, using roots to solve equations is more complicated, because even power functions are not one-to-one.
\begin{definition}[Even Root Property]
For any expression $X$, and real number $c\geq 0$, and even integer $p$, the equation 
$$X^p=c$$
has the same solutions as the two equations
$$X=\sqrt[p]{c}\quad\mbox{or}\quad X=-\sqrt[p]{c}$$
If $c<0$, $X^p=c$ has no solution.
\end{definition}
\begin{note}To simplify our notation, we will sometimes write $X=\pm\sqrt[p]{c}$ as a shortcut of $X=\sqrt[p]{c}\mbox{ or }X=-\sqrt[p]{c}$.\end{note}
\begin{procedure}{Solving Equations with the Even Root Property}
\begin{enumerate*}
\item Isolate the power expression.
\item Apply the even root property to rewrite, if possible.
\item Continue solving.
\end{enumerate*}
\end{procedure}
\begin{example}
Solve each equation below using the Even Root Property.
\begin{lpicvsplit}{2}{7}
\foreach \x/\y/\yy/\eqn in {0/1/5/x^4=16,1/1/7/2(2x+1)^2=50} {
	\regtext{\x*\value{fcols}+\x+\value{fcols}/2}{-\y}{$\eqn$};
	\regtextleft{\x*\value{fcols}+\x}{-\yy}{Check:};
}
\regtext{\value{fcols}/2}{-6}{$(x+3)^6=-1$};

\end{lpicvsplit}
\end{example}
We can solve some equations of this type as polynomial equations.
\begin{example}
Solve $x^4=16$ as a polynomial equation by factoring.
\begin{lpic}{}{5}
\end{lpic}
\end{example}